\section{Discussion}
The arrival schedule determines both which alternatives must be distinguished and how long their distinction can matter. For permanent arrivals with stationary bounded rewards, the schedule-specific minimax cost is an additive spectrum of remaining lifetimes. This result permits very late large batches to be inexpensive, while retaining the possibility that many separated scales each contribute to regret. An all-history causal policy achieves the resulting rate without estimating the future arrival process.

The allocation formulation gives a way to compare schedules before observing rewards. It can distinguish arrival plans with identical endpoints and arm counts, quantify the effect of timing perturbations, and express the marginal contribution of adding candidates. These are comparisons of worst-case learning difficulty. They do not measure the expected benefit of an arm whose mean is known or believed to be high. Using the capacity to design an admission policy would require a model for those benefits and a revised benchmark.

The mechanism experiments qualify how the optimal rate should guide implementation. Clock changes and capacity updates can have opposing finite-sample effects, so a comparison that changes both cannot identify either effect alone. Retaining observations avoids repeated exploration in some environments. The boundary examples and baseline comparisons show that the same minimax rate can conceal substantially different exploration costs near the reward upper bound. The gap bound depends on the smallest positive gaps and may be loose for heterogeneous instances. A general arrival-aware instance complexity, together with a policy attaining it, remains to be established.

The probability lower bound in the supplement shows that spectral difficulty is present with constant probability. The accompanying upper probability bounds come from expectation and are not sharp for small error probabilities. Likewise, the deterministic-time expectation bound alone makes no assertion at reward-dependent stopping times. Reward-dependent arrivals would also invalidate the simple conditioning and deterministic comparator-label arguments used here, and require a separate analysis.

Independent random arrivals change which hard environments can coexist. Theorem~\ref{thm:random-positive-main} resolves this issue for fixed, revealing batch templates: a truncated-lifetime functional replaces the average conditional capacity, and the ratio $D(Q)$ states when averaging is adequate. The result depends on a known endpoint and a fixed template, and its attaining policy uses the schedule information revealed by the first batch. A characterization for nonrevealing random schedules or for a growing collection of batches would have to account for information about future arrivals as well as information about reward means. The simulations illustrate the common-environment restriction without claiming to estimate that more general minimax problem.
